\clearpage
\SourceLeaf{161}\pdfbookmark[0]{Zweites Buch.}{liber-secundus}
\begin{center}
{\Large ZWEITES BUCH.}\par\bigskip
{\large UNTERSUCHUNG DER BAHNEN VON HIMMELSKÖRPERN AUS GEOZENTRISCHEN BEOBACHTUNGEN.}\par\bigskip
\rule{3em}{.3pt}\par\bigskip
\pdfbookmark[1]{Erster Abschnitt.}{liber-secundus-sectio-prima}
{\large ERSTER ABSCHNITT.}\par\medskip
\emph{Bestimmung der Bahn aus drei vollständigen Beobachtungen.}\par\bigskip
\hypertarget{native-161-article115}{}\pdfbookmark[2]{115.}{article-115}115.
\end{center}
Zur vollständigen Bestimmung der Bewegung eines Himmelskörpers in seiner Bahn sind \emph{sieben} Elemente erforderlich. Ihre Zahl vermindert sich jedoch um eins, wenn die Masse des Körpers bekannt ist oder vernachlässigt wird. Bei der Bestimmung einer bisher völlig unbekannten Bahn wird sich die Vernachlässigung der Masse kaum vermeiden lassen, da alle Größen von der Ordnung der Störungen so lange außer Betracht bleiben müssen, bis die Massen, von denen sie abhängen, anderweitig bekannt geworden sind. In der vorliegenden Untersuchung vernachlässigen wir daher die Masse des Körpers und vermindern die Zahl der Elemente auf sechs. Zur Bestimmung einer unbekannten Bahn braucht man demnach ebenso viele Größen, die von den Elementen abhängen und voneinander unabhängig sind. Diese Größen können nur aus den von der Erde aus beobachteten Orten des Himmelskörpers stammen. Da jeder dieser Orte zwei Angaben liefert, nämlich Länge und Breite oder Rektaszension und Deklination, ist es am einfachsten, \emph{drei geozentrische Orte} zu verwenden. Sie genügen im Allgemeinen zur Bestimmung der sechs unbekannten Elemente. Diese Aufgabe ist als die \SourceLeaf{162}wichtigste des vorliegenden Werkes anzusehen und wird daher in diesem Abschnitt mit größter Sorgfalt behandelt.

Fällt jedoch im Sonderfall die Bahnebene mit der Ekliptik zusammen und verschwinden somit alle heliozentrischen wie geozentrischen Breiten ihrer Natur nach, so lassen sich die drei verschwindenden geozentrischen Breiten nicht mehr als drei voneinander unabhängige Angaben betrachten. Die Aufgabe bliebe dann unbestimmt, und unendlich viele Bahnen könnten den drei geozentrischen Orten genügen. In einem solchen Fall müssen also vier geozentrische Längen gegeben sein, damit man die vier übrigen unbekannten Elemente bestimmen kann (die Bahnneigung und die Länge des Knotens fallen weg). Obwohl nach dem Prinzip des Ununterscheidbaren nicht zu erwarten ist, dass ein solcher Fall in der Natur jemals eintritt, liegt die Vermutung nahe, dass die Aufgabe, die für eine genau mit der Ekliptikebene zusammenfallende Bahn gänzlich unbestimmt wird, \emph{bei sehr schwach gegen die Ekliptik geneigten Bahnen} wegen der begrenzten Genauigkeit der Beobachtungen nahezu unbestimmt bleiben muss. Schon die geringsten Beobachtungsfehler können hier die Bestimmung der Unbekannten völlig stören. Um auch diesem Fall Rechnung zu tragen, müssen wir sechs andere Angaben wählen. Zu diesem Zweck werden wir im zweiten Abschnitt zeigen, wie man die unbekannte Bahn aus vier Beobachtungen bestimmt, von denen zwei vollständig sind, während bei den beiden übrigen die Breiten oder Deklinationen fehlen.

Schließlich sind alle unsere Beobachtungen wegen der Unvollkommenheit unserer Instrumente und Sinne nur Annäherungen an die Wahrheit. Eine Bahn, die allein auf den sechs unbedingt erforderlichen Angaben beruht, kann daher noch beträchtliche Fehler aufweisen. Um diese Fehler so weit wie möglich zu vermindern und damit die größtmögliche Genauigkeit zu erreichen, bleibt uns nur, möglichst viele möglichst genaue Beobachtungen zu sammeln und die Elemente so zu verfeinern, dass sie mit allen Beobachtungen möglichst gut übereinstimmen, ohne dass für einzelne Beobachtungen völlige Übereinstimmung verlangt wird. Wie man nach den Grundsätzen der Wahrscheinlichkeitsrechnung eine solche Übereinstimmung erzielt, die überall möglichst eng ist, auch wenn sie nirgends vollkommen ist, werden wir im dritten Abschnitt zeigen.

Auf diese Weise lässt sich die Bestimmung der Bahnen, soweit sich die Himmelskörper darin nach den Gesetzen \textsc{Kepler}s bewegen, bis zu jedem gewünschten Grad der Vollkommenheit führen. Die letzte Verfeinerung lässt sich freilich erst vornehmen, \SourceLeaf{163}wenn auch die Störungen, die die übrigen Planeten in der Bewegung hervorrufen, in die Rechnung einbezogen worden sind. Wie man diese Störungen zu berücksichtigen hat, werden wir im vierten Abschnitt kurz angeben, soweit es für unser Vorhaben von Belang erscheint.
